\section{The distribution protocol of Proposition 7.1}
\label{sec:peps_dyadic_protocol}

The following definition states the clauses
of~\cite[Proposition~7.1]{OpenAI2026PolynomialPEPS} (\texttt{06-geometry.tex},
lines 12--53) as a statement about distributed constructions in the sense of
Definition~\ref{def:peps_distributed_construction} whose chronology is cut into
changes; see~\cite{gap:polypeps_prop71_consequences_interface}. The existence
of such a protocol is Lemma~\ref{prop:peps_dyadic_protocol}.

\begin{definition}[Dyadic distribution protocols and the Section 8 bounds]
    \label{def:peps_dyadic_protocol}
    \lean{TNLean.PEPS.Approximation.DyadicProtocol,
        TNLean.PEPS.Approximation.DyadicProtocol.input,
        TNLean.PEPS.Approximation.DyadicProtocol.trajectory,
        TNLean.PEPS.Approximation.DyadicProtocol.physicalConfigEquiv,
        TNLean.PEPS.Approximation.DyadicProtocol.readout,
        TNLean.PEPS.Approximation.DyadicProtocol.physicalDensity,
        TNLean.PEPS.Approximation.DyadicProtocol.participation,
        TNLean.PEPS.Approximation.HierarchySquare,
        TNLean.PEPS.Approximation.DyadicProtocol.IsDistributionProtocol,
        TNLean.PEPS.Approximation.DyadicProtocol.IsDistributionProtocol.exists_originalCircuit_initialCopies,
        TNLean.PEPS.Approximation.LargeSizeDyadicProtocolConsequences,
        TNLean.PEPS.Approximation.DyadicProtocolBoundsAt,
        TNLean.PEPS.Approximation.LargeSizeDyadicProtocolBounds,
        TNLean.PEPS.Approximation.LargeSizeDyadicProtocolConsequences.largeSizeDyadicProtocolBounds,
        TNLean.PEPS.Approximation.LargeSizeDyadicProtocolBounds.largeSizeDyadicProtocolConsequences,
        TNLean.PEPS.Approximation.DyadicProtocol.IsDistributionProtocol.mono}
    \leanok
    \uses{def:peps_distributed_construction, def:peps_original_circuit,
        def:peps_dyadic_hv_routing}
    A dyadic distribution protocol on the square of side $L$ with local
    dimension $q$ consists of a finite set $P$ of parties, each attached to a
    dyadic square by its scale and block indices, a finite-dimensional private
    memory $H_p$ and a unit vector $x_p\in H_p$ for each party, a number $n$ of
    changes, a register layout $M_t$ of the memory at each stage $t\le n$, and
    for each $t<n$ a contraction circuit $W_t\colon M_t\to M_{t+1}$ whose
    nonprivate gates carry their monomial expansions and their participating
    parties. The first memory is $\bigotimes_pH_p$, and the last consists of one
    physical register of dimension $d_p$ for each party, followed by the private
    registers to be discarded. Each genuine site $v$ has its own party, attached
    to the unit square with address $v$, with $d_p=q$; every other party has
    $d_p=1$. Finally there is a reference vector $r_t\in M_t$ at each stage, with
    $r_0=\bigotimes_px_p$. The protocol assigns a party to each aligned square of
    the padded hierarchy, and records for each change the square of the
    repainting it belongs to, if any, and a list of encoded frames at each stage,
    the encoded copies of $\Omega$. The circuit vector is
    $\Psi_0=r_0$ and
    $\Psi_{t+1}=W_t\Psi_t$.

    For a count constant $C_N$, an expansion constant $C_K$, an error constant
    $C_E$, constants $b$, $k$, $r$ and a unit vector $\Omega$, the protocol
    satisfies the Section~8 bounds if the following hold.
    \begin{enumerate}
        \item Every party is attached to a square of the padded hierarchy,
              there are at most $C_NL^2$ parties, and at most $b$ of them are
              attached to any square. A different party is assigned to each
              aligned square of the padded hierarchy, including the root
              square, at the scale and address of that square. Every other
              party belongs to the repainting of a square of the hierarchy
              other than the root square: it has the scale and address of
              that square, and it participates only in changes of that
              repainting. The changes of the repainting of a square $S$ of
              side $2^j$ form one contiguous segment of the schedule, and
              every party of such a change has scale $j$ or an adjacent
              scale and an address within $b\,2^j$ of the address of $S$ in
              each coordinate. The initial encoded frames are copies of
              $\Omega$ without holes, each held entirely by one party, and
              the initial memory and vector of each party are those of the
              copies it holds; the party holding the first frame holds the
              main copy.
        \item There are at most $C_NL^3$ changes. Each change has at most $b$
              nonprivate gates, each involving at most $b$ parties. Every
              nonprivate gate is a sum of at most $C_KL^k$ monomials with
              absolute coefficient sum at most $C_KL^k$, at most $r$ pair
              effects and at most $r$ pair sources in each monomial. Every
              party participates in at most
              $b$ nonprivate gates over the whole schedule.
        \item Two parties of a nonprivate gate have equal or adjacent scales,
              and their addresses differ in each coordinate by at most $b$
              times the side of the square of either party.
        \item $\lVert W_tr_t-r_{t+1}\rVert\le C_EL^{-20}$ for every $t<n$, and in
              the physical and private coordinates of the last memory
              $r_n=\Omega\otimes\Xi$ with $1-C_EL^{-57}\le\lVert\Xi\rVert\le1$.
              Every intermediate reference vector is the product over the
              encoded frames $F$ of stage $t$, at least one,
              \[
                  r_t=\bigotimes_F\Omega_F,\qquad \Omega_F=K_F\Omega,
                  \qquad \lVert\Omega_F\rVert\le1,\qquad 0<t<n,
              \]
              with the initial frames prepared from $r_0$ by private maps alone and
              the raw registers of the main copy held by the genuine-site parties
              before the final readout change.
    \end{enumerate}
    The existence assertion is that for every $q\ge2$ and $J,\Delta>0$ there are
    constants $C$, $b$, $k$, $r$ and $L_0$, chosen before the size, the
    Hamiltonian and the ground vector, such that for every $L\ge\max\{2,L_0\}$
    and every $(H,E_0,\Omega)\in\mathcal G_{q,J,\Delta}(L)$ there is a protocol
    with these bounds at $C_N=C_K=C_E=C$, the single constant of the source. The
    same requirement at given constants and a given threshold is stated for any
    family of states, consisting for each side $L$ of some vectors on the
    square of side $L$; the existence assertion concerns the family of ground
    vectors of the model. Every
    bound is monotone in each of the three constants, so this assertion is
    equivalent to the one with separate constants $C_N$, $C_K$, $C_E$: given
    those, take $C=\max\{C_N,C_K,C_E\}$. The constants are kept apart because the
    prefactor of the polynomial expansions is much larger than the counts and
    the errors.

    Under the endpoint-coordinate convention just described, these bounds carry
    every clause of~\cite[Proposition~7.1]{OpenAI2026PolynomialPEPS}: strict
    intermediate references are products of encoded frames, whereas the terminal
    reference is expressed in physical and private coordinates as
    $\Omega\otimes\Xi$, with the canonical regrouping of the same auxiliary stack
    proved for the constructed witness rather than retained as a separate field
    of the public protocol predicate. Every intermediate
    reference vector $r_t$, $0<t<n$, has at least one frame
    (\texttt{06-geometry.tex}, lines 38--40), and read per stage, being the
    main copy carries no content beyond the designations at the first stage and
    before the last change. The pair-source bound
    of~\cite[Theorem~5.2]{OpenAI2026PolynomialPEPS}, item~3
    (\texttt{04-compression.tex}, lines 143--147), is part of item~2 above, with
    the same constant $r$ as the pair effects; the recorded monomials contain no
    messages, since messages are expanded before a gate is recorded
    (\texttt{04-compression.tex}, lines 42--45). The gates $W_t$ are the
    selected contractions, as in Section~8 (\texttt{07-assembly.tex}, lines
    29--34). Each nonprivate gate $G$ of $W_t$ equals its recorded expansion,
    $G=\sum_\xi c_\xi M_\xi$, so
    $\lVert G-\sum_\xi c_\xi M_\xi\rVert=0\le L^{-a}$ for every $a$, as
    required of expansions accurate to every fixed inverse power of $L$
    (\texttt{06-geometry.tex}, lines 29--30).
\end{definition}

\begin{lemma}[Distribution through a dyadic hierarchy]
    \label{prop:peps_dyadic_protocol}
    \lean{TNLean.PEPS.Approximation.largeSizeDyadicProtocolConsequences}
    \leanok
    \uses{def:peps_dyadic_protocol}
    For every $q\ge2$ and $J,\Delta>0$ there are constants, uniform in the
    Hamiltonian and in $L$, such that for every sufficiently large $L$ and every
    $(H,E_0,\Omega)\in\mathcal G_{q,J,\Delta}(L)$ there is a protocol that
    distributes $\Omega$ through the padded dyadic hierarchy with all four items
    of~\cite[Proposition~7.1]{OpenAI2026PolynomialPEPS}, applied with patch and
    information power $60$ (\texttt{07-assembly.tex}, line 11). In particular,
    the existence assertion of Definition~\ref{def:peps_dyadic_protocol} holds.
\end{lemma}

\begin{proof}\leanok
    \uses{thm:peps_protocol_from_angular_and_patch_bounds,
        thm:peps_adaptive_patch_constraint}
    The bounds of Definition~\ref{def:peps_dyadic_protocol} are the four items
    of the proposition. Theorem~\ref{thm:peps_protocol_from_angular_and_patch_bounds}
    gives a protocol with these bounds from the adaptive patch constraint of
    Proposition~4.1, which holds by Theorem~\ref{thm:peps_adaptive_patch_constraint}.
\end{proof}

\begin{lemma}[The mesh ratio of the distribution protocol]
    \label{lem:peps_dyadic_protocol_mesh_ratio}
    \lean{TNLean.PEPS.Approximation.protocolMeshRatio,
        TNLean.PEPS.Approximation.protocolHoleRatio_pos,
        TNLean.PEPS.Approximation.protocolMeshRatio_pos,
        TNLean.PEPS.Approximation.protocolMeshRatio_le_pointTreatmentRatio,
        TNLean.PEPS.Approximation.protocolMeshRatio_le_directRepaintingRatio,
        TNLean.PEPS.Approximation.protocolMeshRatio_le_resizeRatio,
        TNLean.PEPS.Approximation.protocolMeshRatio_le_one_sixth,
        TNLean.PEPS.Approximation.protocolMeshRatio_eq}
    \leanok
    \uses{thm:peps_dyadic_block_reading_ratio,
        thm:peps_dyadic_direct_repainting_constants,
        thm:peps_guide_frame_mark_ratio}
    Let $\epsilon_0=1/30000$ and $K_0=20$, and let $\nu$ be the smaller of
    $\nu_2=\min(\epsilon_0,1)/20000$ and
    $\nu_3=\min(\epsilon_0/\max(K_0,1),1)/4$. Then $\nu=1/600000000$, and
    $0<\nu\le\nu_1=\min(\epsilon_0,1)/8$ and $\nu\le1/6$. So $\nu$ serves the
    resizings, the point treatments and the direct repaintings of the protocol
    at once, and satisfies the smallness condition of the constants
    (\texttt{06-geometry.tex}, lines 447--541 and 559--570).
\end{lemma}

\begin{proof}\leanok
    Here $\nu_2=1/600000000$ and $\nu_3=1/2400000$. Since $\nu_2\le\nu_1$ and
    $\nu_2\le1/20000$, both bounds hold.
\end{proof}

\begin{lemma}[The constants of the distribution protocol]
    \label{thm:peps_dyadic_protocol_explicit_constants}
    \lean{TNLean.PEPS.Approximation.largeSizeDyadicProtocol_of_le,
        TNLean.PEPS.Approximation.largeSizeDyadicProtocolBounds,
        TNLean.PEPS.Approximation.protocolHoleRatio,
        TNLean.PEPS.Approximation.rewriteGatePrefactor,
        TNLean.PEPS.Approximation.rewriteGateExponent,
        TNLean.PEPS.Approximation.rewriteGateBound_eq,
        TNLean.PEPS.Approximation.protocolCountConstant,
        TNLean.PEPS.Approximation.protocolLifetimeBound,
        TNLean.PEPS.Approximation.protocolExpansionConstant,
        TNLean.PEPS.Approximation.protocolExpansionExponent,
        TNLean.PEPS.Approximation.protocolErrorConstant,
        TNLean.PEPS.Approximation.protocolPatchBound_protocolMeshRatio}
    \leanok
    \uses{def:peps_dyadic_protocol,lem:peps_dyadic_protocol_mesh_ratio}
    Fix $q\ge2$ and $J,\Delta>0$, let $\epsilon_0=1/30000$, and let
    $\nu=1/600000000$ be the mesh ratio of
    Lemma~\ref{lem:peps_dyadic_protocol_mesh_ratio}. There are an information
    scale $D\ge1$, constants $C,c>0$ with which Proposition~4.1 supplies
    adaptive patches at the radius constant $B_0=20D$ and power $60$, and a
    threshold $L_1$, with the following property. Let
    \begin{align}
        b=\max\{18,\lceil((4+8\epsilon_0+2\nu)/\nu+1)^2\rceil\}
        =2400160003^2=5760768040000960009,
    \end{align}
    $C'=\lceil C\rceil$ and $c'=\lceil c\rceil$, and put
    \begin{align}
        G   & =(C'+1)^{3b}\bigl(16b^2(C'+1)^{6b}+1\bigr)^{2b}, &
        k_K & =3bc'+4b(3bc'+30).
    \end{align}
    For every count constant $C_N\ge5042$, expansion constant $C_K\ge G$ and
    error constant $C_E\ge\max\{10080,2b+1\}$, every $L\ge\max\{2,L_1\}$ and
    every $(H,E_0,\Omega)\in\mathcal G_{q,J,\Delta}(L)$, there is a protocol with
    the Section~8 bounds of Definition~\ref{def:peps_dyadic_protocol} at the
    constants $C_N$, $C_K$, $C_E$, the bound $64\cdot189$ on the lifetime
    participation, the parties per gate and the address distances, the exponent
    $k_K$, and $2b$ pair effects and pair sources per monomial. The prefactor $G$
    of the small-patch rewrites enters only the expansion constant; the
    threshold $L_1$ is that of Proposition~3.4.
\end{lemma}

\begin{proof}\leanok
    \uses{thm:peps_protocol_from_angular_and_patch_bounds,
        thm:peps_adaptive_patch_constraint}
    This is the explicit form of
    Theorem~\ref{thm:peps_protocol_from_angular_and_patch_bounds}, whose
    hypothesis holds by Theorem~\ref{thm:peps_adaptive_patch_constraint}. The
    schedule has at most $5040L^3$ changes, so with the initial and final
    changes the protocol has at most $5042L^3$ changes, and it has fewer than
    $11L^2$ parties. Each frozen auxiliary frame has at most two holes, hence at
    most $10080L^3$ holes in all, and each change has reference error at most
    $(2b+1)L^{-20}$. Each gate is a small-patch rewrite with at most $GL^{k_K}$
    monomials and coefficient sum at most $GL^{k_K}$.
\end{proof}

\begin{lemma}[Consequences of the Section 8 protocol bounds]
    \label{thm:peps_dyadic_protocol_consequences}
    \lean{TNLean.PEPS.Approximation.DyadicProtocol.IsDistributionProtocol.rectangularTraceNorm_physicalDensity_sub_le,
        TNLean.PEPS.Approximation.DyadicProtocol.IsDistributionProtocol.separationBounded,
        TNLean.PEPS.Approximation.DyadicProtocol.participation_eq_card,
        TNLean.PEPS.Approximation.DyadicProtocol.participation_le_of_uses}
    \leanok
    \uses{def:peps_dyadic_protocol, def:peps_dyadic_used_parties}
    Let a protocol satisfy the Section~8 bounds of
    Definition~\ref{def:peps_dyadic_protocol} with $C_E\ge0$, on a square of side
    $L\ge1$, for a unit vector $\Omega$. Then:
    \begin{enumerate}
        \item the physical density of the final circuit vector satisfies
              $\lVert\Tr_{\rm private}|\Psi_n\rangle\langle\Psi_n|
                  -|\Omega\rangle\langle\Omega|\rVert_1\le(2C_NC_E+2C_E)L^{-10}$;
        \item every party network whose links join two parties of a common
              nonprivate gate has link anchors differing in each coordinate by
              at most $b$ times the side of either endpoint square.
    \end{enumerate}
\end{lemma}

\begin{proof}\leanok
    \uses{thm:peps_protocol_density, thm:peps_dyadic_lifetime_parties}
    The changes are contractions, $\lVert r_0\rVert=1$, and the circuit and
    reference vectors agree at the first stage, so the first item is
    Lemma~\ref{thm:peps_protocol_density} with
    $C_{\rm ch}=C_N$ and $C_{\rm step}=C_{\rm hole}=C_E$. The second is the address bound of
    item~3 applied to the two endpoints of a link. For the last, the participation
    of a party is the number of nonprivate gates of all changes in which it
    participates, and Lemma~\ref{thm:peps_dyadic_lifetime_parties} bounds it.
\end{proof}

\begin{lemma}[Retained norm of a protocol with the Section 8 bounds]
    \label{thm:peps_dyadic_protocol_retained_norm}
    \lean{TNLean.PEPS.Approximation.DyadicProtocol.IsDistributionProtocol.retained_norm_and_trace_le}
    \leanok
    \uses{def:peps_dyadic_protocol}
    Let a protocol satisfy the Section~8 bounds of
    Definition~\ref{def:peps_dyadic_protocol} with $C_E\ge0$, on a square of side
    $L\ge1$, for a unit vector $\Omega$. Then the final circuit vector $\Psi_n$
    and its physical density
    $\rho=\Tr_{\rm private}|\Psi_n\rangle\langle\Psi_n|$ satisfy
    \begin{align}
        \lVert\Psi_n\rVert & \ge1-C_EL^{-57}-C_NC_EL^{-17},             &
        1-\tr\rho          & \le2\left(C_EL^{-57}+C_NC_EL^{-17}\right).
    \end{align}
    without renormalizing any intermediate vector
    (\cite[Section~8.1]{OpenAI2026PolynomialPEPS}).
\end{lemma}

\begin{proof}\leanok
    \uses{thm:peps_protocol_retained_norm}
    As in the proof of Lemma~\ref{thm:peps_dyadic_protocol_consequences}, the
    hypotheses of Lemma~\ref{thm:peps_protocol_density} hold with
    $C_{\rm ch}=C_N$ and $C_{\rm step}=C_{\rm hole}=C_E$, and
    Lemma~\ref{thm:peps_protocol_retained_norm} gives both estimates.
\end{proof}
